\par\noindent\textbf{(1)} 把复解析电场写 E⁺(t)=Σ\_m A(t−mT\_R)e\textasciicircum{}\{iω\_c(t−mT\_R)\}e\textasciicircum{}\{imΔφ\_CE\}，其中A为单脉冲包络。平移T\_R后，复场除载波平移外获取固定相位因子；频谱是单脉冲谱与无限周期相位阵列之积。
\par\noindent\textbf{(2)} 在频率f分量上，相邻脉冲干涉相位必须满足 2πfT\_R−Δφ\_CE=2πn，从而 f\_n=n/T\_R+Δφ\_CE/(2πT\_R)≡n f\_rep+f\_CEO。f\_CEO仅模f\_rep有意义，负值可通过重新标记整数n写进标准 [0,f\_rep) 范围。
\par\noindent\textbf{(3)} 自参考时对同一倍频程内低频梳齿f\_n进行倍频得2f\_n=2n f\_rep+2f\_CEO，与高频齿 f\_\{2n\}=2n f\_rep+f\_CEO 拍频，其差为f\_CEO。测得此拍频并反馈可稳定载波包络偏移，重复频率则通过脉冲序列时钟独立锁定。
\par\noindent\textbf{(4)} 腔群速度色散改变T\_R和各频率的相位同步，非线性与增益共同决定稳定锁模；噪声造成频梳齿线宽和相位漂移。Δφ\_CE增加2π不会改变可观测梳齿集合，因此不应把未经取模的偏移频率理解为独立的无限实数参数。
